%%%%%%%%%%%%%%%%%%%%%%%%%%%%%%%%%%%%%%%%%%%%%%%%%%%%%%%%%%%%%%%%%%%%%%%%%%%%%%%
%% Aula exemplo -- vitrine dos recursos da classe guiapratico
%%
%% Este arquivo nao pretende ser uma aula de verdade: ele exercita, um a um, os
%% comandos e ambientes da classe, para servir de consulta rapida na hora de
%% escrever as aulas reais. Copie daqui os trechos que interessarem.
%%
%% ARMADILHA UNICA: \cd e verbatim (\lstinline). Ele NAO pode aparecer dentro
%% do argumento de outra macro -- \objetivos, \campo, \obs, \caption, \section,
%% \item de \objetivos etc. Nesses lugares use \texttt{\textbackslash comando}.
%%%%%%%%%%%%%%%%%%%%%%%%%%%%%%%%%%%%%%%%%%%%%%%%%%%%%%%%%%%%%%%%%%%%%%%%%%%%%%%

% \aula[titulo curto para o sumario]{Titulo completo}
% Abre pagina nova, imprime a faixa cinza "Aula N" com o brasao e zera os
% contadores de secao, figura, tabela, equacao e codigo.
\aula[Recursos do template]{Recursos do Template --- Aula de Exemplo}

% ---------------------------------------------------------------- cabecalho --
% Campos rotulados do topo da aula. Um campo vazio simplesmente nao e impresso.

\objetivos{%
  \item conhecer os comandos de estrutura da classe
        (\texttt{\textbackslash aula}, \texttt{\textbackslash section},
        \texttt{\textbackslash secaoguia}, \texttt{\textbackslash subtituloguia});
  \item exercitar os blocos próprios do guia --- roteiro, tabela de
        preenchimento, figura, listagem de código e citação;
  \item servir de referência rápida durante a escrita das demais aulas.}

\listamaterial{1 computador com \sw{MATLAB}/\sw{Octave} ou \sw{Python};
  1 osciloscópio digital; 1 gerador de funções; 1 protoboard; 1 circuito
  integrado \ci{SN7400}; resistores e capacitores diversos; cabos banana--jacaré.}

% \campo{Rotulo:}{conteudo} -- campo rotulado generico, para quando os campos
% prontos (\objetivos, \listamaterial) nao bastarem.
\campo{Duração:}{2 horas/aula (1h40min).}

% ------------------------------------------------------------------ texto ----
% \introducao imprime o rotulo "Introducao:" (e um \subtituloguia pronto).
\section{Introdução}

Este documento reúne, em um só lugar, tudo o que a classe \cd{guiapratico}
oferece para a escrita de uma aula prática. O texto corrido é composto em
DejaVu Sans, com parágrafos justificados; termos estrangeiros e nomes de
programas passam por \cd|\sw{...}| --- como em \sw{software}, \sw{datasheet} e
\sw{Simulink} --- e códigos de circuitos integrados por \cd|\ci{...}|, que
impede a quebra de linha no meio da sigla: \ci{SN7400}, \ci{CD4017}.

Trechos curtos de código no meio da frase saem com \cd|\cd{...}|: chame
\cd{step(sys)} para obter a resposta ao degrau. Se o trecho contiver \verb|%| ou
chaves desbalanceadas, troque as chaves por outro delimitador, como em
\verb+\cd|a % b|+.

A referência teórica da disciplina é \textcite{ogata2010}; para o projeto de
controladores no domínio da frequência, veja também \cite{nise2017,dorf2018}.
Uma citação com página fica assim: \cite[p.~42]{franklin2019}. Citação de
citação usa \texttt{\textbackslash apud}: \apud{astrom2004}{ogata2010}.

% ATENCAO: \cd nao funciona dentro do argumento de \obs -- use \texttt.
\obs{o comando \texttt{\textbackslash obs} produz este parágrafo iniciado por
\textbf{OBS --}, útil para avisos de segurança ou pegadinhas da montagem.}

% -------------------------------------------------------------- estrutura ----
% Duas familias de titulos convivem no mesmo texto:
%   \section/\subsection/\subsubsection -> numerados (1.1, 1.1.1, 1.1.1.1)
%   \secaoguia/\subtituloguia           -> sem numero, fora do sumario

\section{Títulos numerados}\label{sec:titulos}

Esta é uma \cd{\section}: numerada com o número da aula como prefixo
(\ref{sec:titulos}) e referenciável por \cd|\label|/\cd|\ref|. Por padrão as
seções não aparecem no sumário --- para incluí-las, aumente o \cd{tocdepth} no
preâmbulo do \cd{main.tex}.

\subsection{Subseção}

Um nível abaixo. A numeração acompanha: a Seção~\ref{sec:titulos} gera esta
subseção \thesubsection.

\subsubsection{Sub-subseção}

O nível mais profundo numerado pela classe.

\secaoguia{Título de bloco sem número}

O \cd{\secaoguia} imprime um título em caixa alta, sem numeração e fora do
sumário --- é o mesmo estilo usado em \textbf{REFERÊNCIAS} e nos apêndices.

\subtituloguia{Subtítulo em negrito}

Já o \cd{\subtituloguia} produz um subtítulo indentado, também sem número.

% ---------------------------------------------------------------- formulas ---
\section{Equações}

O \cd{amsmath} já vem carregado pela classe. O contador de equações é zerado a
cada aula (as figuras, tabelas e listagens é que levam o número da aula como
prefixo), e a equação aceita rótulo:

\begin{equation}\label{eq:pid}
  u(t) = K_p\,e(t) + K_i \int_{0}^{t} e(\tau)\,\mathrm{d}\tau
       + K_d \frac{\mathrm{d}e(t)}{\mathrm{d}t}.
\end{equation}

Em Laplace, a \req{eq:pid} vira a função de transferência

\begin{equation}\label{eq:pid-s}
  C(s) = K_p + \frac{K_i}{s} + K_d s
       = K_p\left(1 + \frac{1}{T_i s} + T_d s\right),
\end{equation}

com $T_i = K_p/K_i$ e $T_d = K_d/K_p$. Matemática no meio do texto também
funciona: o sistema de primeira ordem $G(s) = K/(\tau s + 1)$ tem constante de
tempo $\tau$ e ganho estático $K$.

Para citar uma equação no texto use \cd|\req{rotulo}|, que escreve
\enquote{Equação} seguida do número entre parênteses, como na \req{eq:pid-s}.

% ----------------------------------------------------------------- unidades ---
\section{Números e unidades}

A classe carrega o \cd{siunitx}, que é a forma padrão de escrever grandezas
neste guia --- nunca digite a unidade à mão. Ele já vem configurado para o
português: vírgula decimal, milhar separado por espaço fino e faixas escritas
com \enquote{a}.

\begin{itemize}
  \item \cd|\qty{500}{\milli\second}| $\rightarrow$ \qty{500}{\milli\second}
        --- valor com unidade (o espaço correto entre eles é automático e
        não quebra linha);
  \item \cd|\qtyrange{-10,5}{+10,5}{\volt}| $\rightarrow$
        \qtyrange{-10,5}{+10,5}{\volt} --- faixa com a unidade escrita uma
        única vez;
  \item \cd|\unit{\litre\per\minute}| $\rightarrow$ \unit{\litre\per\minute}
        --- só a unidade, para cabeçalhos de tabela e rótulos de eixo;
  \item \cd|\num{-32768}| $\rightarrow$ \num{-32768} --- só o número, com o
        sinal e o agrupamento de milhar corretos;
  \item \cd|\qty{28,3}{\percent}| $\rightarrow$ \qty{28,3}{\percent} ---
        porcentagem (repare que a vírgula decimal é escrita como vírgula).
\end{itemize}

Todos funcionam dentro e fora do modo matemático. Unidades fora do SI usadas
no guia são declaradas na classe com \cd|\DeclareSIUnit|: \cd|\bit|/\cd|\bits|, \cd|\rpm|
e \cd|\voltpp| ($\qty{5}{\voltpp}$, por exemplo).

% ----------------------------------------------------------------- figuras ---
\section{Figuras}

\cd|\figuraguia[largura]{arquivo}{legenda}| coloca a figura centrada, com
legenda numerada, buscando o arquivo em \cd{img/} \emph{ao lado} deste
\cd{.tex} --- não importa de onde o arquivo seja incluído.

\figuraguia[0.8\textwidth]{exemplo.png}{Imagem de exemplo, incluída com
  \texttt{\textbackslash figuraguia}.}

Para uma imagem compartilhada entre várias aulas, guardada em \cd{figuras/} na
raiz do projeto, use \cd|\figuraglobal{figuras/arquivo.pdf}{legenda}|. E para
inserir a imagem sem \sw{float} nem legenda --- dentro de uma tabela, por
exemplo --- existe \cd|\includefigure[largura]{arquivo}|.

% Figura montada "na mao" em vez de \figuraguia, para que o \label venha logo
% depois do \caption -- so assim ela pode ser referenciada por \ref (o
% \figuraguia nao tem onde encaixar o rotulo).
\begin{figure}[H]
  \centering
  \includefigure[0.5\textwidth]{exemplo.png}
  \caption{Circuito usado na prática; aqui o \texttt{\textbackslash label} vem
    logo após o \texttt{\textbackslash caption}, o que permite a referência
    cruzada feita mais adiante, no roteiro.}
  \label{fig:exemplo-ref}
\end{figure}

A chamada no texto é feita com \cd|\rfig{rotulo}|, que escreve
\enquote{Figura} seguida do número --- por exemplo, \rfig{fig:exemplo-ref}.

% ----------------------------------------------------------------- tabelas ---
\section{Tabelas}

O ambiente \cd{tabelaguia} desenha a tabela no estilo do guia: fios pretos e
primeira linha em cinza, negrito. A especificação das colunas segue o
\cd{tabularray} (\cd{c}, \cd{l}, \cd{r}, \cd{X}\ldots).

\begin{table}[H]
  \centering
  \caption{Tabela verdade da porta NAND, preenchida pelo professor.}
  \label{tab:verdade}
  \begin{tabelaguia}{colspec={ccc},width=0.55\textwidth}
    $A$ & $B$ & $S = \overline{A \cdot B}$ \\
    0 & 0 & 1 \\
    0 & 1 & 1 \\
    1 & 0 & 1 \\
    1 & 1 & 0 \\
  \end{tabelaguia}
\end{table}

Nas tabelas de preenchimento à mão, \cd{\vazio} dá altura útil à célula vazia:

\begin{table}[H]
  \centering
  \caption{Medidas do ensaio --- preencher durante a prática.}
  \label{tab:medidas}
  \begin{tabelaguia}{colspec={cccc},width=0.85\textwidth}
    Ensaio & $V_{\text{ent}}$ [\unit{\volt}] & $V_{\text{sai}}$ [\unit{\volt}] & $t_s$ [\unit{\milli\second}] \\
    1 & \vazio & & \\
    2 & \vazio & & \\
    3 & \vazio & & \\
  \end{tabelaguia}
\end{table}

A Tabela~\ref{tab:medidas} também pode ser citada por \cd{\ref}, como a
Tabela~\ref{tab:verdade}.

% ---------------------------------------------------------------- roteiros ---
% \pratica imprime o titulo "Pratica:" alinhado a esquerda.
\pratica

O ambiente \cd{roteiro} numera os passos no formato \textbf{1}~-- e aceita um
nível de subitens (\textbf{1.1}~--).

\begin{roteiro}
  \item Monte na protoboard o circuito da \rfig{fig:exemplo-ref}, com o
        \ci{SN7400} alimentado em \qty{5}{\volt}.

  \item Ajuste o gerador de funções para onda quadrada de \qty{1}{\kilo\hertz},
        amplitude \qty{5}{\voltpp} e nível DC de \qty{2,5}{\volt}:
        \begin{roteiro}
          \item conecte a saída do gerador à entrada $A$;
          \item mantenha a entrada $B$ em nível lógico alto;
          \item verifique a forma de onda no osciloscópio antes de prosseguir.
        \end{roteiro}

  \item Registre na Tabela~\ref{tab:medidas} os valores medidos.

  \item Compare o resultado com a Tabela~\ref{tab:verdade} e justifique as
        divergências, se houver.
\end{roteiro}

% ---------------------------------------------------------------- listagens --
\section{Listagens de código}

Blocos de código têm realce de sintaxe, numeração de linhas e legenda
\cd{Codigo N.M:} --- numerada por aula, como as figuras. Acentuação dentro das
listagens funciona normalmente.

\begin{codigomatlab}[caption={Resposta ao degrau de uma planta de segunda ordem.},
                     label={cod:degrau}]
% Planta de segunda ordem subamortecida
wn   = 2*pi*5;        % frequência natural [rad/s]
zeta = 0.3;           % coeficiente de amortecimento
G    = tf(wn^2, [1 2*zeta*wn wn^2]);

% Malha fechada com controlador proporcional
Kp = 4;
T  = feedback(Kp*G, 1);

step(T); grid on;
title('Resposta ao degrau em malha fechada');
\end{codigomatlab}

O Código~\ref{cod:degrau} usa o ambiente \cd{codigomatlab}. Há ainda
\cd{codigopython}, \cd{codigoc}, \cd{codigocpp}, \cd{codigoarduino},
\cd{codigovhdl}, \cd{codigobash} e \cd{codigotexto} (sem realce). O ambiente
genérico \cd{codigo} aceita \cd{language=} no argumento opcional, junto de
qualquer outra chave do \cd{listings}.

\begin{codigoarduino}[caption={Leitura do sensor e acionamento do atuador.},
                      label={cod:arduino}]
const int PINO_SENSOR  = A0;
const int PINO_ATUADOR = 9;

void setup() {
  Serial.begin(115200);
  pinMode(PINO_ATUADOR, OUTPUT);
}

void loop() {
  int leitura = analogRead(PINO_SENSOR);   // leitura de 0 a 1023
  analogWrite(PINO_ATUADOR, leitura / 4);  // saída PWM de 0 a 255
  Serial.println(leitura);
  delay(10);
}
\end{codigoarduino}

Sem numeração de linhas e sem realce, para saída de terminal:

\begin{codigotexto}[caption={Saída esperada no monitor serial.}, style=semnumeros]
$ python degrau.py
K = 2.00, tau = 3.00 s, theta = 1.50 s
\end{codigotexto}

\subtituloguia{Código em arquivo externo}

Melhor que copiar e colar: o fonte fica em \cd{src/}, roda de verdade e o guia
mostra sempre a versão atual. As chaves \cd{firstline}/\cd{lastline} recortam
um trecho, e \cd{style=semnumeros} desliga a numeração de linhas.

\codigoarquivo[language=Python, firstline=14, lastline=21,
  caption={Identificação pelo método dos dois pontos (arquivo
  \texttt{aulas/aula-exemplo/src/degrau.py}).},
  label={cod:identifica}]{aulas/aula-exemplo/src/degrau.py}

O Código~\ref{cod:identifica} implementa o método de Smith, discutido em
\textcite{astrom2008pid}.

% ------------------------------------------------------------------ listas ---
\section{Listas de apoio}

Além do \cd{roteiro}, valem as listas padrão do \LaTeX, já compactadas pela
classe:

\begin{itemize}
  \item itens sem numeração, para enumerar materiais ou cuidados;
  \item aceitam aninhamento:
    \begin{itemize}
      \item segundo nível;
    \end{itemize}
\end{itemize}

\begin{enumerate}[label=\alph*),leftmargin=1.2cm]
  \item o \cd{enumitem} está carregado, então rótulos personalizados funcionam;
  \item útil para as alíneas exigidas pela ABNT no relatório.
\end{enumerate}

% ---------------------------------------------------------------- pesquisa ---
% \pesquisa imprime o titulo "PESQUISA" centrado.
\pesquisa

Bloco para a parte teórica que o aluno responde em casa.

\begin{roteiro}
  \item Consulte o \sw{datasheet} do \ci{SN7400} \cite{sn7400} e informe a
        corrente máxima de saída em nível baixo ($I_{OL}$).

  \item Explique, com suas palavras, por que o método dos dois pontos usa
        justamente \qty{28,3}{\percent} e \qty{63,2}{\percent} do valor final.

  \item Pesquise na documentação do \sw{Control System Toolbox}
        \cite{mathworks-cst} qual função retorna as margens de ganho e de fase.
\end{roteiro}

\obs{a fonte de toda pesquisa deve ser citada no relatório, conforme a
  \cite{nbr6023} --- veja o Apêndice~B.}
